\section{Treść zadania projektowego}

\textbf{IMSU -- multimedialny inteligentny system uczący:}

Klasyczne kursy internetowe podają materiał każdemu uczniowi w ten sam sposób, choć jedni lepiej uczą się z rysunków i schematów, inni z nagrań i rozmowy, a jeszcze inni przez samodzielne wykonywanie ćwiczeń. Nauczyciel prowadzący zajęcia na żywo dostosowuje się do ucznia odruchowo -- system internetowy musi to robić według jawnych reguł.

Celem projektu jest system, w którym \textbf{nauczyciel} składa w edytorze kurs z lekcji zawierających tekst oraz materiały multimedialne (obrazy, filmy, nagrania, karty PDF, linki), oznaczając każdy materiał formą przekazu: wzrokową (V), słuchową (A) lub kinestetyczną (K). \textbf{Uczeń} po rejestracji wypełnia ankietę VAK, z której system wyznacza jego styl uczenia się oraz wagi poszczególnych form. Na tej podstawie system decyduje, jak podać lekcję: materiały pasujące do stylu ucznia są rozwinięte, pozostałe zwinięte, ale dostępne. Po każdym quizie system aktualizuje wagi form na podstawie wyniku i materiałów, z których uczeń faktycznie korzystał, a gdy inna forma wyraźnie wypada lepiej -- proponuje zmianę stylu wraz z uzasadnieniem. Ostateczną decyzję podejmuje uczeń.

System obsługuje trzy role: ucznia, nauczyciela i administratora, który zarządza kontami użytkowników.